\begin{lemma}
Let $\mathcal H$ be a $3$-uniform hypergraph of maximum degree at most $D$,
and let $e\in E(\mathcal H)$ intersect exactly $3(D-1)$ other edges.  Put
\[
X(e,\mathcal H)=
\left(\bigcup\{h\in E(\mathcal H):h\cap e\ne\emptyset\}\right)\setminus e
\]
and let $\deg_{\mathcal H}(v,x)$ denote the number of edges containing both
$v\in e$ and $x\in X(e,\mathcal H)$.  Then:
\begin{enumerate}
\item every edge $h\ne e$ satisfies $|h\cap e|\le1$;
\item for every $x\in X(e,\mathcal H)$,
$\sum_{v\in e}\deg_{\mathcal H}(v,x)\le D$;
\item for every $v\in e$,
$\sum_{x\in X(e,\mathcal H)}\deg_{\mathcal H}(v,x)=2(D-1)$.
\end{enumerate}
\end{lemma}
\begin{proof}
This is the $k=3$ case of Lemma~\ref{lem: tableop new} of the paper.  By
\noderef{UniformHypergraph}, write $e=\{v_1,v_2,v_3\}$.  The maximum-degree
hypothesis means, by \noderef{MaxDegreeAtMost} and
\noderef{HypergraphDegree}, that each $v_i$ lies in at most $D$ edge indices,
hence at most $D-1$ edge indices other than $e$.

If an edge $h\ne e$ contained two vertices of $e$, then the incidence count
of neighboring edges by the vertex of $e$ they contain would count $h$ twice.
The total incidence count is at most $3(D-1)$, so the number of distinct
neighbors of $e$ would be at most $3(D-1)-1$, contradicting the hypothesis
that there are exactly $3(D-1)$ neighbors.  This proves the first item.

Fix $x\in X(e,\mathcal H)$.  By the first item, no edge can contain $x$ and
two distinct vertices of $e$.  Hence the sets of edge indices counted by the
three numbers $\deg_{\mathcal H}(v_i,x)$ are disjoint subsets of the set of
all edge indices containing $x$.  Their total size is therefore at most
$\deg_{\mathcal H}(x)\le D$, proving the column-sum bound.

Finally fix $v_i\in e$.  The first item makes the three classes of neighbors,
classified by their unique vertex in $e$, disjoint.  Since their union has
size $3(D-1)$ and each class has size at most $D-1$, the class containing
$v_i$ has exactly $D-1$ edges.  Each such edge has three vertices by
\noderef{UniformHypergraph}, meets $e$ only in $v_i$, and therefore
contributes exactly two incidences to
$\sum_x\deg_{\mathcal H}(v_i,x)$.  The row sum is $2(D-1)$.
\end{proof}
